% Pursuing Stacks (cotes 134-2 to 134-8): archive cote 134-2, site page 10 (PDF page 11).
% provisional contents. Draft compared: none.
% AI-transcribed by Claude Opus 5.5 (claude-opus-5-5, adaptive thinking, effort high), Message Batches API, batch ps1 (msgbatch_01Rv5hCaNTtUt4JHx1yqoHPL), 2026-11-12; not yet checked by eye.
\scanpage{134-2/10}
\hand{100. Homology versus cohomology}
\hand{101. Abelian integration $Lf_!^{ab}$ (versus abelian cointegra-%}
\hand{\textcircled{102} Integrators are abelianizators.} \hand{tion $Rf_*$ )}
\hand{103. Integrators versus cointegrators}
\hand{104. Overall review on abelianization (1): case of} \inserted{pseudo-topoi}
\hand{105. Review (2): duality equivalences for ``algebraic''}
\hand{topoi\unsure{,|.} \deleted{and abelian topoi}.}
\hand{106. Review (3): A formulaire for the basic integra-%}
\hand{tion and cointegration operations $*$ and Hom.}
\hand{107. Review (4): Case of general ground ring $k$.}
\hand{108. Review (5): Homology and cohomology}
\hand{(absolute case)}
\hand{\textcircled{109.} Review (6): A further step in linearization:}
\hand{coalgebra structures $P\to P\otimes_k P$ in (Cat).}
% NOTE: a handwritten horizontal line runs across the page here
\margin{\fbox{VI} Schematization.}
\hand{\textcircled{110.} More wishful thinking on ``schematization'' of homotopy types.}
\hand{\textcircled{111.} Complexes of ``unipotent bundles'' as models, and}
\hand{``schematic'' linearization.}
\hand{112. Postnikov dévissage and Kan condition for}
\hand{schematic complexes.}
\hand{113. ``Soft'' versus ``hard'' Postnikov dévissage; $\pi_1$ as}
\hand{a group scheme.}
\hand{114. Outline of a program.}
\hand{115. $L(X)$ as the pro-quasicoherent substitute for $\underline{O}_k^{(X)}$.}
\hand{116. The need and the drive.}
\hand{117. ``Schematic'' versus ``formal'' homology and}
\hand{cohomology invariants.}
\hand{\textcircled{118.} The homotopy groups $\pi_i$ as derived functors}
\hand{of the ``Lie functor''. Lack of satisfactory}
\hand{models for $S^2$ and $S^3$.}

\mspage{\hand{10}}
\hand{-- -- --}
\hand{-- -- -- (à suivre)}
